\chapter{Diseño}

\section{Visión general}

El sistema se diseñó como un pipeline de generación científica y software. Su
entrada es una descripción en lenguaje natural de un problema de toma de
decisiones. Su salida es un conjunto de modelos Python autocontenidos,
registrados y probados, que pueden cargarse posteriormente en el laboratorio
virtual. Entre ambos extremos se generan artefactos intermedios: informes de
investigación, formulaciones matemáticas, especificaciones JSON, código,
pruebas, trazas y memorias.

La arquitectura no delega todo el proceso en un único modelo de lenguaje. En su
lugar, separa el sistema en cuatro capas principales: orquestación, implementada
por el Router; agentes especializados; artefactos persistentes; y memoria y
recuperación de conocimiento.

\begin{figure}[H]
\centering
\begin{verbatim}
Usuario / CLI / Web / Evaluaciones
        |
        v
   +---------+
   | Router  |
   | estados |
   | review  |
   | trazas  |
   +----+----+
        |
  +-----+-------------+----------------+
  |                   |                |
  v                   v                v
Agentes          Artefactos          Memoria
R/F/R/B          MinIO + PG          Neo4j + Qdrant + PG
subagentes       estado/trazas       recuperación
herramientas     modelos             ciclo de vida
\end{verbatim}
\caption{Arquitectura lógica del sistema.}
\end{figure}

Esta separación responde a una necesidad práctica: las salidas de los modelos
de lenguaje son útiles, pero deben estar acotadas por contratos, herramientas,
validaciones, revisión humana y persistencia. El diseño busca que cada etapa
produzca algo inspeccionable y que los errores puedan localizarse en el nivel
correcto.

\section{Problema de diseño}

El problema no es solo generar código, sino transformar conocimiento científico
en comportamiento simulado. Esa transformación tiene varios cambios de
representación:

\begin{figure}[H]
\centering
\begin{verbatim}
Problema en lenguaje natural
        |
        v
Paradigmas cientificos relevantes
        |
        v
Formulaciones matematicas
        |
        v
Especificaciones adaptadas al entorno
        |
        v
Modelos Python + pruebas
\end{verbatim}
\caption{Transformación conceptual del pipeline.}
\end{figure}

Cada cambio de representación reduce ambigüedad. La investigación decide qué
marcos científicos son relevantes. La formalización convierte esos marcos en
variables, parámetros y ecuaciones. El razonamiento adapta las ecuaciones a un
entorno concreto, con percepciones y acciones disponibles. La construcción
traduce la especificación a código ejecutable y pruebas. Si todo esto lo hiciera
un único agente, sería difícil saber si un fallo procede de una mala fuente
bibliográfica, una formulación incoherente, una mala adaptación al entorno o un
bug en Python.

\section{Router y máquina de estados}

El Router es el coordinador central. Se implementa como una máquina de estados
explícita, no como una planificación libre de un LLM. Esta decisión permite
reanudar ejecuciones, registrar eventos, controlar revisiones humanas y colocar
la escritura de memoria solo después de que una salida haya sido aceptada.

\begin{figure}[H]
\centering
\begin{verbatim}
CLASSIFY_UMBRELLA
  -> RESEARCH -> REVIEW_RESEARCH -> MEMORY_RESEARCH
  -> FORMALIZE -> REVIEW_FORMALIZE -> MEMORY_FORMALIZE
  -> GET_ENV_SPEC
  -> REASON -> REVIEW_REASON -> MEMORY_REASON
  -> BUILD -> REVIEW_BUILD -> MEMORY_BUILD
  -> DONE
\end{verbatim}
\caption{Máquina de estados simplificada.}
\end{figure}

Las etapas \texttt{MEMORY\_*} son condicionales: se ejecutan cuando la
infraestructura de memoria está disponible. Si Neo4j, Qdrant o los embeddings no
están configurados, el pipeline puede continuar en modo degradado. En cambio,
PostgreSQL y MinIO son servicios centrales, porque almacenan ejecuciones,
metadatos, artefactos y estado.

El Router mantiene un \texttt{PipelineState} con la etapa actual, los paradigmas
aprobados, las formulaciones seleccionadas y las especificaciones aprobadas.
Además, emite trazas de agentes, herramientas y artefactos. Esto permite
reconstruir qué ocurrió en una ejecución concreta y continuar desde el estado
persistido si el proceso se interrumpe.

\section{Por qué hay cuatro agentes principales}

El sistema tiene cuatro agentes principales de generación:
\textit{Researcher}, \textit{Formalizer}, \textit{Reasoner} y
\textit{Builder}. También existen componentes auxiliares, como el Classifier,
el DeepResearcher, los subagentes y el MemoryAgent. La división principal en
cuatro agentes se debe a que cada uno resuelve un tipo de problema distinto.

\begin{table}[H]
\begin{center}
\begin{tabular}{|p{2.6cm}|p{4.2cm}|p{5.6cm}|} \hline
\textbf{Agente} & \textbf{Pregunta que resuelve} & \textbf{Razón de separación} \\ \hline
Researcher & Qué paradigmas científicos explican el problema. & Requiere búsqueda, lectura de literatura y síntesis científica. \\ \hline
Formalizer & Cómo se expresan esos paradigmas con variables, parámetros y ecuaciones. & Requiere razonamiento matemático y comparación entre formulaciones alternativas. \\ \hline
Reasoner & Cómo se implementa una formulación en un entorno concreto. & Requiere mapear conceptos abstractos a percepciones, acciones y recompensas reales. \\ \hline
Builder & Cómo se traduce una especificación a código probado. & Requiere escribir Python, respetar contratos, generar tests y corregir fallos ejecutables. \\ \hline
\end{tabular}
\caption{Motivo de la división en cuatro agentes principales.}
\end{center}
\end{table}

Esta división permite que la revisión humana intervenga con precisión. Si la
investigación omite un paradigma, se corrige en la etapa de investigación. Si la
formulación matemática es mala, se rerunea la formalización. Si la adaptación al
entorno usa acciones inexistentes, se corrige el Reasoner. Si el código falla,
se corrige el Builder. Un agente monolítico mezclaría estos niveles y haría más
difícil diagnosticar los errores.

\section{Workflow de agentes}

\subsection{Classifier}

Antes de lanzar la investigación completa, el sistema ejecuta un clasificador
ligero. Su función es intentar asociar el problema a un paradigma canónico ya
conocido. Esto evita que el sistema fragmente la memoria creando nombres
distintos para conceptos equivalentes. Por ejemplo, una ejecución sobre
\textit{temporal-difference learning} puede reutilizar un paraguas más amplio
como \textit{reinforcement learning} si ya existe en el grafo.

El Classifier no produce un artefacto principal. Su salida se usa como señal
interna para orientar al Researcher y para reducir duplicados de identidad en la
memoria.

\subsection{Researcher y DeepResearcher}

El Researcher descubre paradigmas de toma de decisiones relevantes para el
problema. Primero consulta la memoria para obtener paradigmas conocidos y
entidades relacionadas. Después usa búsqueda web y subagentes de investigación
profunda para cubrir huecos reales. Finalmente sintetiza un informe global y
emite una lista estructurada de paradigmas.

El DeepResearcher es un subagente especializado en un único paradigma. Produce
un informe profundo con fundamentos, postulados, supuestos, predicciones,
localización primaria, conceptos clave, variables identificadas y referencias.
El Researcher padre recibe un resumen, pero el informe completo se guarda como
artefacto para las etapas siguientes.

Los artefactos principales de esta etapa son:

\begin{verbatim}
research/{run_id}/report.md
research/{run_id}/deep/{paradigm}.md
\end{verbatim}

Después de esta etapa, el humano revisa qué paradigmas continúan. También puede
pedir que se investigue un paradigma adicional si considera que falta una línea
científica importante.

\subsection{Formalizer}

El Formalizer recibe los paradigmas aprobados y genera formulaciones
matemáticas. No trabaja sobre todos los paradigmas en un único contexto: lanza
un \texttt{FormalizerSubAgent} por paradigma. Cada subagente lee el informe
profundo correspondiente y escribe un archivo de formulaciones.

Cada formulación debe incluir variables, parámetros, ecuaciones y lógica de
decisión. Además, las formulaciones de un mismo paradigma deben ser
estructuralmente distintas. No basta con renombrar parámetros; el sistema busca
alternativas como control algebraico, ecuaciones diferenciales, modelos
probabilísticos, aprendizaje por refuerzo o políticas por umbrales.

El artefacto principal es:

\begin{verbatim}
research/{run_id}/formulations/{paradigm}.md
\end{verbatim}

Tras esta etapa, el humano selecciona qué formulaciones merecen continuar. Esta
selección es importante porque no todas las variantes generadas son igual de
útiles o implementables.

\subsection{Reasoner}

El Reasoner adapta las formulaciones seleccionadas al entorno de simulación. Su
entrada no es solo la formulación matemática, sino también el archivo
\texttt{env\_spec.json}, que define el espacio de percepción, las acciones
disponibles, los recursos y las recompensas.

Cada \texttt{ReasonerSubAgent} lee:

\begin{verbatim}
deep/{paradigm}.md
formulations/{paradigm}.md
env_spec.json
\end{verbatim}

y produce una especificación JSON por formulación válida:

\begin{verbatim}
models/{run_id}/reasoner/{paradigm}/{formulation}.json
\end{verbatim}

La especificación JSON es el contrato entre Reasoner y Builder. Contiene
variables, parámetros, reglas de actualización, lógica de decisión, mapeo al
entorno, comportamientos esperados y referencias. Antes de escribir una
especificación, el Reasoner valida que las variables estén definidas, que no
haya dependencias circulares incoherentes, que las acciones existan en el
entorno y que las percepciones usadas sean observables o computables.

Si una formulación no es válida, el Reasoner no fuerza una implementación:
escribe un informe JSON con \texttt{status: "invalid"} y una lista de problemas.
Así el fallo queda documentado en el nivel correcto.

\subsection{Builder}

El Builder traduce especificaciones JSON aprobadas a modelos Python. Cada
\texttt{BuilderSubAgent} trabaja con una especificación concreta y genera:

\begin{verbatim}
models/{run_id}/builder/{paradigm}/{formulation}_model.py
models/{run_id}/builder/{paradigm}/test_{formulation}.py
\end{verbatim}

El modelo generado debe ser autocontenido y cumplir el contrato de modelo
esperado. El Builder también genera pruebas y ejecuta
\texttt{uv run pytest}. Si las pruebas fallan, el subagente puede corregir el
código y rerunear los tests hasta un límite de intentos.

Además, el Builder valida implementabilidad antes de escribir código. Si la
lógica de decisión es ambigua, si usa claves de percepción inexistentes o si los
comportamientos esperados no son testeables, guarda un informe de validación en
lugar de producir código defectuoso.

\section{Uso de subagentes}

Los subagentes se usan en las etapas donde existen unidades de trabajo
independientes. Un paradigma puede formalizarse sin esperar a otro; una
formulación puede adaptarse sin mezclar su contexto con otras; una
especificación puede implementarse sin cargar todo el historial del pipeline.

El diseño con subagentes aporta cuatro ventajas:

\begin{itemize}
\item \textbf{Contexto más pequeño:} cada subagente recibe solo los artefactos
que necesita.
\item \textbf{Paralelismo:} varios paradigmas o especificaciones pueden
procesarse con \texttt{asyncio.gather}.
\item \textbf{Aislamiento de fallos:} si una formulación falla, no invalida las
demás.
\item \textbf{Artefactos estables:} cada subagente escribe en rutas
deterministas, fáciles de revisar y reutilizar.
\end{itemize}

Esta es una diferencia importante frente a un único agente grande. El objetivo
no es solo reducir tiempo, sino mantener límites claros entre artefactos y
evitar que errores de una rama contaminen el resto de la ejecución.

\section{Herramientas de los agentes}

Los agentes no actúan únicamente con texto. Cada uno dispone de herramientas
limitadas a su responsabilidad. Esta restricción reduce el espacio de error:
un agente de investigación puede buscar, pero no escribir código; el Builder
puede ejecutar tests, pero no decidir qué paradigmas son científicamente
relevantes.

\begin{table}[H]
\begin{center}
\small
\begin{tabular}{|p{2.8cm}|p{4.4cm}|p{5.2cm}|} \hline
\textbf{Agente} & \textbf{Herramientas} & \textbf{Motivo} \\ \hline
Researcher & \texttt{retrieve\_knowledge}, \texttt{web\_search}, \texttt{launch\_deep\_research}, \texttt{read\_report} & Reutilizar paradigmas conocidos, buscar huecos reales, lanzar investigación profunda y leer informes completos antes de sintetizar. \\ \hline
Deep researcher & \texttt{retrieve\_knowledge}, \texttt{search\_papers}, \texttt{web\_search} & Combinar memoria previa, búsqueda académica con DOI/autores y contexto web general. \\ \hline
Formalizer sub. & \texttt{read\_file}, \texttt{write\_file}, \texttt{retrieve\_knowledge} & Leer informes profundos, escribir formulaciones y recuperar patrones matemáticos ya validados. \\ \hline
Reasoner sub. & \texttt{read\_file}, \texttt{write\_file}, \texttt{retrieve\_knowledge} & Leer investigación, formulaciones y entorno; escribir specs JSON; reutilizar rangos de parámetros y mapeos de entorno. \\ \hline
Builder sub. & \texttt{read\_file}, \texttt{write\_file}, \texttt{run\_tests}, \texttt{retrieve\_knowledge} & Leer specs, escribir modelo y tests, ejecutar validación real y recuperar patrones de implementación previos. \\ \hline
MemoryAgent & extracción estructurada, escritura en grafo, indexado vectorial, resolución de conflictos & No es un agente conversacional abierto; ejecuta una tubería fija para persistir conocimiento aceptado. \\ \hline
\end{tabular}
\caption{Herramientas disponibles por agente.}
\end{center}
\end{table}

La herramienta transversal más importante es \texttt{retrieve\_knowledge}. Solo
está disponible cuando la infraestructura de memoria lo permite. Su función no
es sustituir la búsqueda web, sino introducir conocimiento acumulado de
ejecuciones anteriores: paradigmas ya conocidos, variables, postulados, rangos
de parámetros, patrones de mapeo al entorno y errores de implementación
detectados previamente.

\section{Human-in-the-loop}

La revisión humana forma parte del diseño del pipeline, no es una operación
externa. El Router llama a un \texttt{FeedbackPort}, que abstrae tres modos de
uso: revisión por CLI, revisión por WebSocket y autoaprobación para el entorno
de evaluación. Así, la lógica del Router no depende de si el usuario revisa en
terminal, interfaz web o pruebas automáticas.

\begin{table}[H]
\begin{center}
\begin{tabular}{|p{3cm}|p{4.8cm}|p{4.8cm}|} \hline
\textbf{Punto de revisión} & \textbf{Decisión humana} & \textbf{Por qué existe} \\ \hline
Investigación & Aprobar paradigmas o pedir investigar uno adicional. & Evita que el pipeline continúe con marcos científicos irrelevantes o incompletos. \\ \hline
Formalización & Seleccionar paradigmas y formulaciones concretas. & Reduce variantes matemáticas a las que tienen sentido para implementar. \\ \hline
Entorno & Proporcionar \texttt{env\_spec.json}. & El sistema necesita una definición autoritativa de acciones, percepciones y recompensas. \\ \hline
Razonamiento & Aprobar specs JSON, rechazarlas o rerunear formalización. & Distingue errores de adaptación al entorno de errores matemáticos previos. \\ \hline
Construcción & Aprobar modelos generados, rechazar código o rerunear Reasoner. & Separa problemas de implementación de problemas de especificación. \\ \hline
\end{tabular}
\caption{Puntos de revisión humana del pipeline.}
\end{center}
\end{table}

El aspecto clave es que la revisión no solo decide si continuar o parar. También
puede redirigir el pipeline al nivel adecuado. Por ejemplo, un error conceptual
en una formulación debe volver al Formalizer; una mala conversión de esa
formulación al entorno debe volver al Reasoner; un bug de Python debe volver al
Builder.

\section{Artefactos y persistencia}

El sistema es \textit{storage-first}: los artefactos reales se guardan en MinIO
y sus metadatos se registran en PostgreSQL. Esto evita depender de carpetas
locales y permite que CLI, servidor web y evaluaciones compartan el mismo modelo
de persistencia.

\begin{verbatim}
research/{run_id}/
  report.md
  deep/{paradigm}.md
  formulations/{paradigm}.md
  env_spec.json
  pipeline_state.json
  trace.jsonl

models/{run_id}/
  reasoner/{paradigm}/{formulation}.json
  builder/{paradigm}/{formulation}_model.py
  builder/{paradigm}/test_{formulation}.py
  builder/{paradigm}/{formulation}_validation.json
\end{verbatim}

MinIO almacena el contenido de los artefactos. PostgreSQL almacena ejecuciones,
metadatos, modelos aprobados y filas de memoria. Esta separación evita guardar
grandes bloques de texto o código en SQL, pero conserva referencias
relacionales para listar ejecuciones, registrar modelos y recuperar estado.

La traza \texttt{trace.jsonl} registra eventos de agentes, herramientas y
artefactos. Junto con \texttt{pipeline\_state.json}, permite auditar una
ejecución y reanudarla sin reconstruir el contexto desde cero.

\section{Contrato de modelo generado}

La salida final no hereda de una clase concreta del simulador. Usa duck typing:

\begin{lstlisting}[language=Python]
def decide(self, perception: dict) -> Action
def update(self, action, reward, new_perception) -> None
def get_state(self) -> dict
\end{lstlisting}

Esta decisión reduce el acoplamiento entre la fase de generación y la fase de
simulación. El modelo generado puede cargarse dinámicamente siempre que cumpla
la interfaz esperada.

La separación entre \texttt{decide} y \texttt{update} es crítica. El método
\texttt{decide} debe ser de solo lectura: elige una acción usando el estado
actual y la percepción. El método \texttt{update} recibe la recompensa y la nueva
percepción, y es el único que puede modificar el estado interno. Esta regla
evita que el modelo aprenda o cambie variables antes de que el entorno haya
aplicado realmente la acción.

Además, \texttt{get\_state} debe exponer variables internas y una clave
\texttt{q\_values}. Esta información permite a la fase de simulación visualizar
alternativas de acción, confianza relativa y evolución de puntuaciones internas.

\section{Sistema de memoria y conocimiento}

\subsection{Objetivo de la memoria}

La memoria convierte el pipeline en un sistema acumulativo. Sin memoria, cada
ejecución empieza desde cero. Con memoria, los resultados aceptados se
transforman en conocimiento reutilizable: paradigmas, variables, postulados,
formulaciones, parámetros, modelos, relaciones y hechos textuales.

La regla central es:

\begin{verbatim}
salida de agente
  -> revision humana
  -> artefacto aceptado
  -> memoria
\end{verbatim}

El sistema no guarda cualquier borrador generado por un LLM. Solo escribe en la
memoria después de la revisión humana. Esta decisión reduce ruido, evita
persistir formulaciones rechazadas y hace que la memoria sea explicable: lo que
hay en ella procede de artefactos aceptados.

\subsection{MemoryAgent}

El MemoryAgent no es un agente libre que decide qué hacer en cada turno. Es una
tubería determinista que usa modelos de lenguaje para tareas concretas, pero
mantiene fijo el flujo de control:

\begin{figure}[H]
\centering
\begin{verbatim}
Artefacto aceptado
      |
      v
Extraccion estructurada
      |
      v
Canonicalizacion y normalizacion
      |
      v
Escritura en Neo4j + revision del grafo
      |
      v
Resolucion de duplicados/conflictos
      |
      v
PostgreSQL + Qdrant dense/sparse
\end{verbatim}
\caption{Flujo interno del MemoryAgent.}
\end{figure}

La extracción convierte el artefacto aceptado en tres tipos de salida:

\begin{itemize}
\item \textbf{nodos:} entidades como paradigmas, variables, papers,
postulados, formulaciones o modelos;
\item \textbf{relaciones:} enlaces de grafo como \texttt{SUPPORTS},
\texttt{IMPLEMENTS}, \texttt{USES\_VARIABLE}, \texttt{DERIVES\_FROM} o
\texttt{HAS\_PARAMETER};
\item \textbf{hechos:} frases atómicas que pueden buscarse y recuperarse en
ejecuciones futuras.
\end{itemize}

Antes de escribir, el sistema normaliza identidades, poda nodos sin relación,
materializa relaciones estructurales y resuelve propuestas \texttt{\_\_NEW\_\_}
contra paradigmas ya existentes. Esto evita que la base de conocimiento se
llene de duplicados semánticos.

\subsection{Qué se guarda en cada almacén}

La memoria usa varios almacenes porque cada uno resuelve un problema distinto.

\begin{table}[H]
\begin{center}
\begin{tabular}{|p{2.4cm}|p{9.8cm}|} \hline
\textbf{Almacén} & \textbf{Responsabilidad} \\ \hline
MinIO & Artefactos completos: informes, formulaciones, specs, código, tests, estado y trazas. \\ \hline
PostgreSQL & Ejecuciones, metadatos, modelos registrados y ciclo de vida de memorias: confianza, validez temporal, accesos y supersesión. \\ \hline
Neo4j & Identidad y topología del conocimiento: nodos científicos y relaciones entre ellos. \\ \hline
Qdrant & Recuperación de hechos mediante dos colecciones: \texttt{memories\_dense} para embeddings y \texttt{memories\_sparse} para BM25. \\ \hline
\end{tabular}
\caption{Responsabilidad de cada almacén de persistencia.}
\end{center}
\end{table}

PostgreSQL es la fuente de verdad para confianza y validez temporal. Neo4j
guarda relaciones, pero no decide si una memoria está vigente. Qdrant permite
buscar, pero sus payloads no son la fuente de verdad de la confianza. MinIO
conserva el texto completo de los artefactos originales.

\subsection{Tipos de memoria}

Las memorias de fase 1 se guardan en la tabla
\texttt{pipeline\_memories}. Cada fila incluye contenido, \texttt{namespace},
\texttt{memory\_type}, etapa de origen, ejecución, importancia, confianza,
contadores de corroboración y contradicción, accesos, intervalo de validez y
posible memoria sucesora.

\begin{table}[H]
\begin{center}
\small
\begin{tabular}{|p{2.5cm}|p{2.8cm}|p{2.7cm}|p{4.3cm}|} \hline
\textbf{Etapa} & \textbf{Namespace} & \textbf{Tipo} & \textbf{Para qué sirve} \\ \hline
Researcher & \texttt{paradigm} & \texttt{semantic} & Recordar paradigmas, postulados, autores, variables y conceptos científicos. \\ \hline
Formalizer & \texttt{formulation} & \texttt{semantic} & Reutilizar patrones matemáticos, ecuaciones, variables y alternativas de modelado. \\ \hline
Reasoner & \texttt{formulation} & \texttt{semantic} & Recordar mapeos entre formulaciones, percepciones, acciones, parámetros y comportamientos esperados. \\ \hline
Builder & \texttt{model} & \texttt{procedural} & Reutilizar patrones de implementación, tests, errores de código y decisiones del contrato \texttt{DecisionModel}. \\ \hline
Memoria & \texttt{meta} & \texttt{episodic} / \texttt{reflection} & Registrar contradicciones, consolidaciones y observaciones sobre la propia memoria. \\ \hline
\end{tabular}
\caption{Tipos de memoria de la fase 1.}
\end{center}
\end{table}

Esta separación permite búsquedas más precisas. Una consulta del Researcher
puede restringirse a \texttt{paradigm}; una consulta del Builder puede buscar en
\texttt{model}; una consolidación puede trabajar sobre memorias \texttt{meta}.

\subsection{Ciclo de vida de las memorias}

La memoria no es un log append-only. Cada hecho tiene ciclo de vida. Cuando se
extrae un hecho, el sistema calcula un hash normalizado para evitar duplicados
exactos dentro de la misma ejecución y etapa. Después busca hechos similares en
Qdrant. Si encuentra candidatos, clasifica la relación:

\begin{itemize}
\item \textbf{duplicado:} no se crea una memoria nueva;
\item \textbf{corroboración:} aumenta la confianza de la memoria existente;
\item \textbf{enriquecimiento:} se crea una versión nueva y la anterior queda
supersedida;
\item \textbf{contradicción:} baja la confianza de la memoria anterior, se
crea una nueva memoria y se añade una memoria meta sobre el conflicto.
\end{itemize}

La confianza evoluciona de forma explícita. Una corroboración suma confianza,
una contradicción la reduce, el acceso mediante recuperación la incrementa
ligeramente y las memorias antiguas pueden decaer si no se usan. El objetivo es
que el sistema no solo recuerde, sino que pueda cambiar de criterio cuando
aparecen evidencias mejores.

\subsection{Recuperación de conocimiento}

La herramienta \texttt{retrieve\_knowledge} combina varias señales:

\begin{figure}[H]
\centering
\begin{verbatim}
Consulta del agente
      |
      v
Reescritura de consulta
      |
      +--> Qdrant denso
      +--> Qdrant BM25
      +--> Neo4j local 2-hop
      |
      v
Fusion RRF + reranking
      |
      v
CRAG / fallback web si hay baja confianza
      |
      v
Contexto inyectado al agente
\end{verbatim}
\caption{Flujo de recuperación de conocimiento.}
\end{figure}

La búsqueda densa captura similitud semántica. BM25 captura coincidencias
exactas de nombres, símbolos, autores o identificadores. El grafo permite
recuperar relaciones que no son necesariamente cercanas en texto, como que una
formulación implementa un paradigma o que una variable pertenece a un postulado.
Después se fusionan los resultados y se rerankean. Si la confianza no es
suficiente, el sistema aplica una revisión correctiva y puede complementar con
búsqueda web.

Esto hace que la memoria sea útil pero no ciega. El sistema no inyecta
automáticamente cualquier resultado almacenado: pondera confianza, actualidad,
relevancia y señales de corrección.

\section{Degradación controlada}

El sistema distingue servicios obligatorios y opcionales. PostgreSQL y MinIO son
obligatorios para ejecutar el pipeline de forma persistente. Neo4j, Qdrant,
embeddings y reranking son opcionales desde el punto de vista de la generación:
si faltan, se pierde recuperación y memoria avanzada, pero el pipeline puede
seguir investigando, formalizando, razonando y construyendo modelos.

\begin{table}[H]
\begin{center}
\small
\begin{tabular}{|p{3cm}|p{9.5cm}|} \hline
\textbf{Servicio} & \textbf{Comportamiento si no está disponible} \\ \hline
PostgreSQL & No se puede mantener correctamente la persistencia central de ejecuciones, modelos y memorias. \\ \hline
MinIO & No se pueden guardar artefactos, estado ni trazas de forma estable. \\ \hline
Neo4j & Se pierde recuperación relacional y escritura de grafo, pero pueden existir artefactos y modelos. \\ \hline
Qdrant & Se pierde búsqueda densa/BM25 sobre hechos de memoria. \\ \hline
Embeddings & La recuperación semántica y el reordenado se degradan o se desactivan. \\ \hline
\end{tabular}
\caption{Degradación ante fallo de servicios.}
\end{center}
\end{table}

\section{Decisiones de diseño principales}

Las decisiones más importantes del diseño son:

\begin{itemize}
\item separar investigación, formalización, razonamiento y construcción;
\item usar un Router determinista en lugar de un agente planificador global;
\item usar revisión humana antes de escribir memoria;
\item guardar artefactos persistentes en MinIO y metadatos en PostgreSQL;
\item usar Neo4j para relaciones, Qdrant para recuperación y PostgreSQL para
confianza y validez temporal;
\item mantener el contrato \texttt{DecisionModel} por duck typing;
\item validar antes de implementar y ejecutar pruebas después de generar código;
\item permitir degradación controlada cuando falta infraestructura opcional.
\end{itemize}

En conjunto, estas decisiones convierten el sistema en algo más que una cadena
de prompts. El resultado es una arquitectura de agentes con límites definidos,
artefactos revisables, memoria acumulativa, validación ejecutable y capacidad de
reanudación.
